\chapter{CLIO: Constraint-Leveraged Inference and Optimization}
\label{app:clio}

\noindent \CLIO --- Constraint-Leveraged Inference and Optimization --- is the
recursive constraint repair framework developed as the self-correction layer of
the Yarncrawler architecture. This appendix distinguishes the framework from the
convergent system of Cheng, Broadbent, and Chappell 	tep{cheng2025clio}, develops the formal
specification, and identifies the key design differences.

\section{Two Versions of CLIO}

The acronym \CLIO has been used independently in two frameworks with overlapping
but distinct concerns. The Cheng--Broadbent--Chappell system (hereafter
CBC-CLIO) is an inference-time orchestration layer over large language models
that uses recursive confidence signals to steer reasoning without post-training.
The framework developed here (hereafter RC-CLIO for Recursive Constraint repair)
addresses a different problem: maintaining coherence across overlapping theoretical
frameworks in a distributed conceptual manifold.

The convergences are real. Both systems use recursive depth as a control variable.
Both treat uncertainty reduction as a positive signal. Both expose their reasoning
process rather than hiding it behind post-hoc rationalization. The CBC-CLIO
finding that correct reasoning exhibits a negative uncertainty gradient over time
is consistent with the RC-CLIO corrigibility condition: a system whose
reconstruction errors are being corrected shows decreasing unresolved residue.

The divergences are equally real. CBC-CLIO optimizes toward known correct answers
on benchmarks. RC-CLIO operates in settings where the admissibility conditions
themselves are being constructed rather than given---where there is no external
oracle to verify the output. CBC-CLIO's uncertainty signal is self-reported by
the model. RC-CLIO's uncertainty signal is the non-trivial cohomological residue
of the theoretical manifold, detectable through gluing failures at framework
boundaries.

\section{The RC-CLIO Specification}

\begin{definition}{RC-CLIO Repair Cycle}
A single \emph{RC-CLIO repair cycle} over a theoretical manifold
$(\mathcal{M}, \{(U_i, \varphi_i)\})$ consists of four phases:

\textbf{Phase 1: Seam detection.} Compute the seam penalty
$L_{\rm seam}(x) = \sum_{i < j} w_i(x)w_j(x)\|f_i(x) - f_j(x)\|^2$
over all overlapping charts. Identify regions where $L_{\rm seam}$ exceeds
a threshold $\epsilon_{\rm repair}$.

\textbf{Phase 2: Cocycle identification.} For each high-seam region, identify
the Čech 1-cocycle responsible: the mismatch between restrictions of adjacent
local sections to their shared overlap.

\textbf{Phase 3: Repair morphism construction.} Introduce new morphisms into
the affected local sections that trivialize the cocycle. The repair morphism
must satisfy locality (modifying only sections near the tear), type safety
(remaining within the admissible semantic category), and stigmergic monotonicity
(reducing the seam penalty without increasing it elsewhere).

\textbf{Phase 4: Residue annotation.} Record the repair in the system's
provenance structure: which seam was repaired, by which morphism, and what
alternative repairs were available but not applied.
\end{definition}

\section{The Corrigibility Condition in RC-CLIO}

\begin{definition}{RC-CLIO Corrigibility}
An RC-CLIO system is \emph{corrigible} with respect to a target manifold
$\mathcal{M}^*$ if:
\begin{enumerate}
  \item The seam detection phase correctly identifies all regions where the
    current manifold $\mathcal{M}$ differs from $\mathcal{M}^*$.
  \item The repair morphisms introduced in Phase 3 reduce $d(\mathcal{M},
    \mathcal{M}^*)$ monotonically.
  \item The residue annotation in Phase 4 is complete: all applied morphisms
    are recorded with their alternatives.
\end{enumerate}
A system fails corrigibility if it introduces repair morphisms that reduce the
seam penalty $L_{\rm seam}$ while increasing $d(\mathcal{M}, \mathcal{M}^*)$:
this is proxy stabilization at the level of the repair process itself.
\end{definition}

\section{Relationship to CBC-CLIO}

The CBC-CLIO system implements a version of Phase 1 (seam detection through
uncertainty monitoring) and Phase 3 (repair through recursive re-invocation
of the reasoning process). It does not implement Phase 2 (cohomological
identification of the obstruction) or Phase 4 (provenance annotation of
repairs). This is appropriate given CBC-CLIO's operational context: it is
designed for single-session scientific question-answering, not for long-horizon
theoretical manifold maintenance.

\claimC The CBC-CLIO uncertainty oscillation finding---that correct answers
show negative uncertainty gradients while incorrect answers show positive or
oscillating gradients---is consistent with Phase 1 of RC-CLIO. In RC-CLIO
terms: a system whose seam penalties are converging to zero is on a corrigible
trajectory; a system whose seam penalties are oscillating has entered a repair
cycle that is not reducing $d(\mathcal{M}, \mathcal{M}^*)$ despite reducing
$L_{\rm seam}$ locally. The oscillation is the signature of proxy stabilization
in the repair process.
